\bookchapter{제24장: Root 권한과 BL 잠금 심층 분석}

\chapterintro{이 장은 제20장 내용의 심화 확장으로, 스마트폰의 기저 메커니즘을 진정으로 이해하거나, 플래싱을 시도하려는 사용자에게 적합합니다. 참고로, 이 작업에는 일정 위험이 있으므로 신중하게 진행해 주세요.}

\booksection{하나, Bootloader 심층 이해}

Bootloader는 스마트폰이 시작될 때 첫 번째로 실행되는 프로그램으로, 그 임무는 하드웨어 초기화, 운영체제 커널 로드입니다. Bootloader를 스마트폰의 “문지기”로 이해할 수 있습니다.

\booksubsection{BL 잠금의 작동 원리:}

\begin{itemize}
  \item \textbf{서명 검증}: BL 잠금은 부팅 이미지의 디지털 서명을 검증하며, 제조사 서명된 시스템만 부팅할 수 있습니다.
  \item \textbf{데이터 보호}: 새 시스템을 플래싱해도, 원래 계정 비밀번호 없이는 사용자 데이터에 접근할 수 없습니다.
\end{itemize}

\booksubsection{하드웨어 수준 저장과 검증}

BL 잠금의 잠금 상태(Locked/Unlocked) 및 제조사 공개 키는 일반적으로 보호된 하드웨어 영역(예: Qualcomm의 devinfo 파티션, eFuse 또는 RPMB 파티션)에 저장되어, 소프트웨어 변조를 방지합니다. 매 부팅 시, Bootloader는 이 안전 영역에서 상태를 읽고, 그에 따라 엄격한 서명 검증을 수행할지 결정합니다.

\booksubsection{디지털 서명 검증}

제조사는 시스템 업데이트나 이미지를 출시할 때 자신의 개인키로 서명합니다. 기기 부팅 시 Bootloader는 내장된 공개키를 사용하여 시스템 이미지(예: boot.img, recovery.img)를 검증합니다. 작동 방식은 다음과 같습니다:

\begin{enumerate}
  \item 시스템 이미지에 대해 해시 값을 다시 계산합니다.
  \item 공개키를 사용하여 이미지에 포함된 서명을 복호화하여 해시 값을 얻습니다.
  \item 이 두 해시 값을 비교하여 일치하면 검증을 통과합니다. 이를 통해 시스템 이미지가 제3자에 의해 변조되지 않았음을 보장합니다.
\end{enumerate}

\booksubsection{보안 부팅 체인}

BL 잠금은 전체 보안 부팅 체인의 첫 번째 고리입니다. 검증이 통과된 후에야 다음 단계의 부트로더(예: 커널)를 로드하고 실행하여 시스템 부팅 환경의 신뢰성을 단계별로 보장합니다.

\booksubsection{잠금 vs 잠금 해제: BL 잠금의 양면적 영향}

BL 잠금의 장단점을 더 직관적으로 이해하기 위해 5가지 핵심 차원에서의 차이를 대비 표로 보여드립니다:

\begin{center}
\begin{tabularx}{\linewidth}{|X|X|X|}
\hline
\textbf{핵심 역할} & \textbf{잠금 상태(기본값)} & \textbf{잠금 해제 후 상태} \tabularnewline
\hline
시스템보안 & 서명되지 않은 비공식 시스템 플래싱을 방지하고 시스템 무결성과 저수준 보안을 보호\newline 높은 보안: 공식 서명 이미지만 로드하여 악성 펌웨어 변조를 효과적으로 방어 & 보안 저하: 타사 ROM 플래싱 허용, 시스템 무결성 검증 메커니즘 무효화, 악성 소프트웨어가 Root 권한을 얻으면 시스템을 깊이 변조할 수 있음 \tabularnewline
시스템 안정성 & 스마트폰이 제조사가 테스트하고 최적화한 공식 시스템에서 실행되어 안정적이고 원활한 작동을 보장\newline 안정적: 공식 펌웨어 사용, 기능 완전, OTA 업그레이드 신뢰성 & 위험 증가: 타사 ROM은 호환성 문제가 있을 수 있으며, 시스템 충돌, 벽돌, 하드웨어 기능 이상(예: 카메라, 급속 충전)을 유발할 수 있고, OTA 업데이트가 실패할 수 있음 \tabularnewline
보증 혜택 & 제조사는 BL 잠금을 통해 안정적인 공식 시스템을 제공하여 보증 서비스를 제공\newline 보증 유효: 공식 무료 보증 서비스 이용 가능(인적 손상 외) & 보증 무효: 대부분의 제조사는 BL 잠금 해제를 자발적 보증 포기로 간주하며, 재잠금하더라도 과거 잠금 해제 기록이 감지되어 보증이 거부될 수 있음 \tabularnewline
데이터보안 & 잠금 해제 과정에서는 일반적으로 데이터 강제 삭제가 이루어져 구 데이터가 새 시스템 환경에서 악용되는 것을 방지\newline 데이터 보안: 일상 사용에서 BL 잠금은 악성 프로그램의 시스템 파티션 수정 능력을 제한 & 데이터 위험: 잠금 해제 시 모든 사용자 데이터 강제 삭제; 잠금 해제 후 올바르게 암호화하지 않으면 데이터 저장이 안전하지 않을 수 있음; 분실 후 플래싱하여 데이터를 더 쉽게 추출할 수 있음 \tabularnewline
기능 제한 & 보안 환경에 의존하는 기능(예: 결제, 은행 앱)의 정상 작동을 보장\newline 기능 완전: 모든 공식 기능 정상 사용 가능 & 기능 제한: 일부 보안 민감 앱(예: Google Pay, 은행 앱)은 BL 잠금 해제 상태를 감지하고 실행을 거부할 수 있음; DRM 서비스가 저하되어 HD 영상 재생에 영향 \tabularnewline
\hline
\end{tabularx}
\end{center}

\booksection{BL 잠금 해제의 상세 절차(샤오미를 예로)}

\begin{enumerate}
  \item \textbf{잠금 해제 권한 신청}: 샤오미 잠금 해제 도구 웹사이트에서 잠금 해제를 신청합니다.
  \item \textbf{계정 연결}: 스마트폰 설정에서 샤오미 계정을 연결합니다.
  \item \textbf{대기 기간}: 계정 등급에 따라 7-30일 대기 필요.
  \item \textbf{잠금 해제 작업}: 샤오미 잠금 해제 도구를 사용하여 컴퓨터에 연결하여 잠금 해제.
\end{enumerate}

\booksection{Root의 다양한 방안}

\begin{itemize}
  \item \textbf{Magisk}: 현대적인 Root 방안, Systemless 설계, Root 상태 숨김 가능.
  \item \textbf{Kernelsu}: 커널 기반 Root 방안, 더 저수준, 은폐성 우수.
  \item \textbf{SuperSU}: 오래된 Root 방안, 현재 업데이트 중지됨.
\end{itemize}

\booksection{플래싱 입문}

플래싱이란 스마트폰에 새 운영체제를 설치하는 것을 말합니다. 일반적인 플래싱 유형:

\begin{itemize}
  \item \textbf{공식 ROM 플래싱}: 공식 시스템 설치, 벽돌 복구 또는 다운그레이드용.
  \item \textbf{타사 Recovery 플래싱}: TWRP 등, 백업 복구, 타사 ROM 플래싱용.
  \item \textbf{타사 ROM 플래싱}: LineageOS, Pixel Experience 등, 다른 시스템 경험.
\end{itemize}

\booksection{삼성 Knox: BL 잠금 해제의 비가역적 대가}

다른 안드로이드 제조사와 달리, 삼성 스마트폰의 BL 잠금 해제 메커니즘에는 고유한 “퓨즈” 디자인이 있습니다:\textbf{삼성 스마트폰의 경우, 잠금 해제 시 Knox 퓨즈가 발동되어 일부 보안 기능(예: Samsung Pay)을 영구적으로 잃고, 보호된 모든 데이터가 삭제되며, 보증에 영향을 미칩니다}. 구체적으로:

\begin{itemize}
  \item \textbf{퓨즈는 일회용}: 스마트폰 내부에 eFuse가 있어, BL 잠금 해제가 감지되면 영구적으로 끊어져 복구 불가.
  \item \textbf{영향받는 기능}: Samsung Pay(교통카드/은행카드), Knox 보안 폴더, 기업 MDM 관리, 일부 건강보험 앱이 모두 작동 불능.
  \item \textbf{강제 데이터 삭제}: 퓨즈 발동 후 스마트폰이 강제 공장 초기화되며, 모든 로컬 데이터가 삭제됨(클라우드 데이터는 영향 없음).
  \item \textbf{보증 영향}: 공식적으로 Knox 퓨즈 발동 후 보증 서비스를 더 이상 제공하지 않음.
\end{itemize}

\begin{tcbwarn}{신중하게 생각하기}
삼성의 Knox 퓨즈는\textbf{비가역적}으로, 한 번 발동되면 공장 상태로 돌아갈 수 없습니다. 단지 Root가 필요한 일부 기능만 사용하고 싶다면, Root가 필요 없는 대체 방안이 있는지 먼저 확인하는 것이 좋습니다.

\end{tcbwarn}

\textbf{중요 경고}\par

플래싱에는 위험이 있으며, 신중하게 작업해야 합니다. 다음을 유발할 수 있습니다:

\begin{itemize}
  \item 스마트폰 벽돌(부팅 불가)
  \item 보증 상실
  \item 데이터 손실
  \item 결제 기능 이상
\end{itemize}

플래싱 전에 반드시 전체 백업을 수행하고, 튜토리얼을 주의 깊게 읽으며, 기종과 버전이 일치하는지 확인하세요.

\begin{tcbinfo}{추천 리소스}
\begin{itemize}
  \item \textbf{XDA Developers}: 세계 최대의 안드로이드 개발자 커뮤니티, 풍부한 플래싱 튜토리얼과 ROM 리소스 제공.
  \item \textbf{Coolapk(쿠안)}: 국내 유명한 디지털 커뮤니티, 많은 중국어 튜토리얼과 토론 제공.
  \item \textbf{각 기종 포럼}: 샤오미 커뮤니티, OnePlus 커뮤니티 등, 공식 및 민간 튜토리얼 제공.
\end{itemize}

\end{tcbinfo}

\booksection{자주 묻는 질문 (Q\& A)}

\textbf{Q1: BL 잠금 해제 후 복구 가능한가요?}\par

\textbf{A:}재잠금 가능합니다. 일반적으로 공식 전체 라인즈 패키지를 플래싱하고 특정 플래싱 스크립트(예: 샤오미의\texttt{flash\_all\_lock.\allowbreak bat})를 사용하여 재잠금할 수 있습니다. 단, 재잠금 후에도 잠금 해제로 인해 트리거된 일부 보안 퓨즈(예: 삼성의 Knox)는 되돌릴 수 없을 수 있습니다.

\textbf{Q2: 잠금 해제 후 Root를 얻는 방법은?}\par

\textbf{A:}BL 잠금 해제는 첫 번째 단계일 뿐입니다. Root는 일반적으로 다음이 필요합니다:

\begin{enumerate}
  \item 커스텀 Recovery 플래싱(예: TWRP).
  \item Recovery를 통해 Magisk 등 Root 관리자를 플래싱하여 boot 이미지를 패치하여 Root 권한을 얻음.
\end{enumerate}
